\documentclass[10pt,a4paper]{amsart}
\usepackage[margin=19mm]{geometry}
\usepackage{amsmath,amssymb,mathtools}
\usepackage[scr=rsfs,cal=euler]{mathalfa}
\usepackage{xfrac}
\usepackage[unicode,hidelinks,bookmarks=false]{hyperref}
\newcommand{\ie}{i.e.\ }
\newcommand{\cf}{cf.\ }
\newcommand{\resp}{respectively\ }
\newcommand{\Subsection}{\subsection}
\providecommand{\nobreakdash}{\nobreak\hbox{-}}
\setlength{\emergencystretch}{2em}
\multlinegap0pt
\usepackage{amssymb}
\newtheorem{lemma}{Lemma}
\def\R{{\mathfrak R}}
\title{Uniqueness and nondegeneracy of positive solutions to elliptic equations with Robin boundary conditions}
\author{Mengyao Chen, Massimo Grossi, Qi Li}
\date{Source: arXiv:2607.29461v1 (CC BY 4.0)}
\begin{document}
\maketitle

\noindent\emph{Source and licence.} Uniqueness and nondegeneracy of positive solutions to elliptic equations with Robin boundary conditions. Version arXiv:2607.29461v1 (CC BY 4.0), distributed by arXiv under the Creative Commons Attribution 4.0 International licence, http://creativecommons.org/licenses/by/4.0/, as recorded in the arXiv licence metadata for this exact version and shown on the abstract page https://arxiv.org/abs/2607.29461v1. Full licence notice: \texttt{licenses/arxiv-2607.29461-CC-BY-4.0.txt}.\par\smallskip

\noindent\textit{Editorial scope.} Counted unit: the article's lemma lemma3.1 (source0.tex lines 744--746). The article's complete original proof of it is delivered with it (source0.tex lines 749--877, one \texttt{proof} environment of 129 lines). No mathematics is written, repaired or re-derived: the statement and the proof are the article's own lines, in the article's own notation, and no retained line is substituted or re-flowed.  Retained uncounted as the source-local context the proof needs: lines 165--171; lines 292--297.  Nothing was omitted from any retained span, so the package carries no omission record.  No retained line contains a citation, so the wrapper carries no bibliography and no bibliography entry.  No retained line contains a figure environment, an image-inclusion command or drawing code, so no image is delivered and the package declares no asset dependency.  Editorial formatting only, in addition: an AMS \texttt{amsart} preamble that restates the article's own declarations and macros used by the retained lines, namely the environments \texttt{lemma} on separate counters, together with the standard packages those lines need.  The retained environments print in this document as Lemma 1 in order of appearance, whereas the article prints the same items with the numbering of its own full document.  The complete native source of the article as distributed in the arXiv e-print for this exact version is bundled unchanged as \texttt{sources/206488/source0.tex}.
\begin{equation}\label{eq0.1}
\begin{cases}
-\Delta u=u^p,\quad &\text{in}\ \Omega,\\
u>0, &\text{in}\ \Omega,\\
\frac{\partial u}{\partial \nu}+\beta u=0, &\text{on}\ \partial \Omega,
\end{cases}
\end{equation}

\begin{equation}\label{1.6}
\begin{cases}
-\Delta h-pu^{p-1}h=0, \quad &\text{in}\ \Omega,\\
\frac{\partial h}{\partial \nu}+\beta h=0 &\text{on}\ \partial\Omega,
\end{cases}
\end{equation}

\begin{lemma}\label{lemma3.1}
Suppose that $\beta>0$, $p$ is subcritical, and $\Omega$ is a ball. If $N=2$, or $N\ge 3$ with $\beta\ge N-2$, then each positive radial solution of \eqref{eq0.1} is nondegenerate in the space of radial functions.
\end{lemma}

\begin{proof}
Without loss of generality, we assume that $\Omega=B_1$. Let $u$ be a positive radial solution of \eqref{eq0.1}. It is sufficient to prove that \eqref{1.6} has only the trivial solution $h\equiv0$ in $H^1_{rad}(\Omega)$. It follows from \eqref{eq0.1} and \eqref{1.6} that 
$$
(p-1)\int_{\Omega}u^phdx=\int_{\Omega}[(-\Delta h)u-(-\Delta u)h]dx=0.
$$
Since $p>1$, we obtain 
\begin{equation}\label{eq3.6}
\int_{\Omega}u^phdx=0.
\end{equation}
If $h\equiv 0$, the proof is complete. Otherwise, \eqref{eq3.6} implies that $h$ must change sign. 

Define the nodal set of $h$ by
$$
\mathcal{M}=\overline{\{x\in\Omega : h(x)=0\}}.
$$
%In the following, we first prove 
%$$
%\mathcal{M}\cap \partial \Omega=\emptyset.
%$$
%If not, we may suppose that the nodal line intersects with the boundary \partial\Omega$. Let y and z be two points which the nodal set of \mathcal{M} intersects with \partial\Omega$, and \Omega_0\subset{\Omega}$ be a subdomain constructed by the boundary \partial\Omega$ and the nodal line from y to z. Then the boundary of \Omega_0$ is defined by 
%$$
%\Gamma_1=\bar{\Omega}_0\cap \partial\Omega,\quad \Gamma_2=\bar{\Omega}_0\cap \mathcal{M}.
%$$
%Note that h has a fixed sigh in \Omega_0$, we may suppose that h>0$. Since 
%\begin{equation}\label{eq3.7}
%-\Delta h=pu^{p-1}h>0,\quad \text{in}\ \Omega_0,
%\end{equation}
%then by Hopf's Lemma we get \frac{\partial h}{\partial\nu}<0 on the boundary \Gamma_2$. We observe that h=0 on \Gamma_2$, then \nabla h is parallel to \nu of \Gamma_2$. For x\in\Gamma_2$, it follows \frac{\partial h}{\partial\nu}<0 that \nabla h=-|\nabla h|\nu$. Since  y,z\in \Omega and \Omega=B_1$, then 
%$$
%\nabla h(y)\cdot y=-\beta h(y)=0,\quad \nabla h(z)\cdot z=-\beta h(z)=0,
%$$
%which implies that \nabla h is tangent to \Gamma_1$ at y and z. Moreover, by the fact that \nabla h=-|\nabla h|\nu on \Gamma_1$, we have that the direction of \nabla h is inward at y and z. Note that 
%$$
%\frac{\partial h}{\partial \nu}(x)=\nabla h\cdot x=-\beta h(x)<0,\quad \forall x\in\Gamma_1\backslash\{y,z\},
%$$
%we can see that the direction of \nabla h is also inward. 
% 
Since $h$ is radial, we claim that $\mathcal{M}\cap \partial\Omega=\emptyset$. Indeed, if there exists $x_0 \in \mathcal{M} \cap \partial\Omega$, we have $h(1)=0$. By the Robin boundary condition, this implies $h'(1)=-\beta h(1)=0$. Since $h$ satisfies the ODE
$$\begin{cases}
-h''-\frac{N-1}rh'=pu^{p-1}(r)h&\text{in }(0,1)\\
h(1)=h'(1)=0
\end{cases}
$$
uniqueness for the Cauchy problem immediately yields $h \equiv 0$, which is a contradiction.
 


%Let 
%$$
%r_0=\sup\{r : h(r)\neq 0 \}.
%$$
%Then r_0\le 1. If r_0<1, we have that  h(r)=0 for all r_0<r\le 1. Moreover, we get h^\prime(r_0+0)=0. On the other hand, without loss of generality, suppose that h(r)>0 in (r_0-\delta,r_0) for some small \delta>0. Since -\Delta h=pu^{p-1}h>0 in (r_0-\delta,r_0), then by Hopf's lemma we get h^{\prime}(r_0-0)<0. Thus, h^\prime(r_0) does not exist since h^{\prime}(r_0-0)<h^{\prime}(r_0+0). This is a contradiction. Thus, we derive that r_0=1. This implies that h has a fixed sign near the boundary \partial\Omega$. Using Hopf's lemma again, we can deduce that h^\prime(1) has a fixed sign.  This is a contradiction with h^\prime(1)=0. 
%}


Define 
$$
w=x\cdot\nabla u.
$$
Since $u$ is radial, then we get 
$$
w=x\cdot \nabla u=ru^\prime(r).
$$
A direct computation yields 
$$
\Delta w=2\Delta u+x\cdot \nabla (\Delta u)=-2u^p-pu^{p-1}w.
$$
%For a fixed point y\in\R^N$, let  T be a differential operator defined by
%$$
%T=(x_1-y_1)\frac{\partial}{\partial x_1}+\cdots+(x_N-y_N)\frac{\partial}{\partial x_N}.
%$$
%Then 
%\begin{equation}\label{3.6}
%\Delta T=T\Delta +2\Delta.
%\end{equation}
%Define 
%$$
%w=Tu,
%$$
%by \eqref{3.6} and \eqref{eq0.1} we get
%$$
%\Delta w=2\Delta u+(x-y)\cdot\nabla(\Delta u)=-2u^p-pu^{p-1}w.
%$$
This implies that $w$ satisfies 
\begin{equation}\label{3.7}
-\Delta w-pu^{p-1}w=2u^p.
\end{equation}
Multiplying \eqref{3.7} by $h$, we get
\begin{equation}\label{3.8}
\int_{\Omega}\nabla w\nabla hdx-\int_{\partial\Omega}\frac{\partial w}{\partial\nu}hdS-\int_{\Omega}pu^{p-1}whdx=2\int_{\Omega}u^phdx.
\end{equation}
On the other hand, multiplying \eqref{1.6} by $w$, we derive 
\begin{equation}\label{3.9}
\int_{\Omega}\nabla h\nabla wdx+\int_{\partial\Omega}\beta hwdS-\int_{\Omega}pu^{p-1}hwdx=0.
\end{equation}
It follows from \eqref{eq3.6}, \eqref{3.8} and \eqref{3.9} that 
\begin{equation}\label{3.10}
\int_{\partial\Omega}\Big(\frac{\partial w}{\partial\nu}+\beta w\Big)hdS=-2\int_{\Omega}u^phdx=0.
\end{equation}
Observe that 
$$
\nabla w=\nabla (ru^\prime(r))=(ru^\prime)^\prime\frac{x}{r}=(u^\prime+ru^{\prime\prime})\frac{x}{r},
$$
then for any $x\in\partial\Omega$ we have 
$$
\frac{\partial w}{\partial\nu}+\beta w=\nabla w\cdot x+\beta w=u^{\prime\prime}(1)+(\beta+1)u^\prime(1).
$$
Since $h$ has a fixed sign on $\partial\Omega$, by \eqref{3.10} we obtain  
\begin{equation*}
u^{\prime\prime}(1)+(\beta+1)u^\prime(1)=0.
\end{equation*}
Moreover, by the Robin boundary condition, we get $u'(1) = -\beta u(1)$, which yields 
\begin{equation}\label{eq3.11}
u^{\prime\prime}(1)=-(\beta+1)u^\prime(1)=\beta(\beta+1)u(1).
\end{equation}
On the other hand, it follows from \eqref{eq0.1} that 
$$
-u^{\prime\prime}(r)-\frac{N-1}{r}u^\prime=u^p,\quad 0<r\le 1.
$$
This implies that 
\begin{equation}\label{e3.12}
u^{\prime\prime}(1)=-u^p(1)-(N-1)u^\prime(1)=[(N-1)\beta-u^{p-1}(1)]u(1).
\end{equation}
Thus, it follows from \eqref{eq3.11} and \eqref{e3.12} that 
$$
u^{p-1}(1)=(N-2-\beta)\beta\le0.
$$
This is a contradiction since $u(1)>0$. 
\end{proof}

\end{document}
